\section{Empirical Signatures}\label{sec:empirical_signatures}

The experiments target qualitative signatures of the validity field:
an interior optimum over memory scale, roughness-dependent horizon movement, a
real-stream analogue, positive separation between validity loss and alarm,
regime-dependent operational control, and an explicit conformal-calibration
readout. The reported sweeps use repeated Monte Carlo runs and are meant to
identify shape, ordering, and sign rather than to estimate constants precisely.

\subsection{Interior Memory Optimum}\label{sec:empirical_memory_optimum}

The Gaussian horizon sweep isolates the variance--staleness trade-off directly.
Small windows are variance dominated, large windows are stale, and the minimum
lies at an interior horizon.

\begin{figure}[!htbp]
\centering
\safeincludegraphics[width=0.88\linewidth]{artifacts/figures/gaussian/fig_ucurve}
\caption{Interior risk optimum over memory length. Tracking error first falls
with horizon and then rises once staleness overtakes the remaining variance
gain. The optimum remains finite across drift strengths.}
\label{fig:memory_optimum}
\end{figure}

This interior optimum is the basic empirical signature of validity under drift.
Under stationarity, more memory would continue to help. Under drift, the horizon
that reduces variance also ages the evidence. Around the minimizing horizon,
the curve displays a flat near-optimal region, which is the empirical analogue
of the useful-memory region defined from the envelope $\Phi$.

\subsection{Roughness-Dependent Horizon Scaling}
\label{sec:empirical_roughness_scaling}

The roughness-family experiment varies the roughness index directly. For each
H\"older-type regime, the surrogate error obeys the same structural
decomposition as the theory,
\[
  \mathcal E(n) \approx \sigma n^{-1/2} + \zeta n^H,
\]
and the experiment traces the empirically optimal horizon across a sweep of
$\zeta$ values.

\begin{figure}[!htbp]
\centering
\safeincludegraphics[width=0.9\linewidth]{artifacts/figures/roughness_family/fig_roughness_family}
\caption{Roughness-dependent horizon scaling. Each path class exhibits a finite
optimum, and the location of that optimum changes with $H$. The horizon is a
function of drift regularity, not a fixed window rule.}
\label{fig:roughness_family}
\end{figure}

The empirical optima separate by roughness class, matching the path-dependent
horizon logic of \cref{cor:temporal_validity_horizon}.

\subsection{Real-Stream Memory Diagnostic}

A similar signature appears on a real nonstationary stream. On ELEC2
\citep{Harries1999ELEC2}, the diagnostic compares a short anchor block with the
preceding memory window in $W_1$, yielding a real-data analogue of the
finite-memory U-curve.

\begin{figure}[!htbp]
\centering
\safeincludegraphics[width=0.88\linewidth]{artifacts/figures/elec2/fig_elec2_ucurve}
\caption{ELEC2 real-stream diagnostic on the \texttt{nswprice} marginal. The
mean anchor-block $W_1$ error exhibits a finite interior optimum and a
near-optimal useful-memory band, matching the validity signature on a real
stream.}
\label{fig:elec2_ucurve}
\end{figure}

The optimum occurs at a finite memory scale with a short near-optimal band. The
same retention pattern therefore appears on a real stream: very short windows
are too noisy, while very long windows become stale relative to the present
block. On the robustness grid, the \texttt{nswprice} marginal retains an
interior optimum in all tested configurations, and the \texttt{nswdemand}
marginal shows the same interior pattern with more movement in the optimizing
scale; see \Cref{tab:elec2_robustness}.

\subsection{Validity Loss Before Detection}\label{sec:empirical_validity_lag}

Temporal validity and changepoint evidence are different clocks. To make that
separation explicit, consider a stream whose local drift rate ramps upward over
time. The horizon $n_t^*$ contracts as the ramp steepens, while a long
operating horizon remains fixed. Define
\[
  \Delta_{\mathrm{val-det}} \coloneqq t_{\mathrm{detect}} - t_{\mathrm{valid}},
\]
where $t_{\mathrm{valid}}$ is the first time the operating horizon persistently
exceeds the validity horizon, and $t_{\mathrm{detect}}$ is the first
detector alarm that follows.

\begin{figure}[!htbp]
\centering
\safeincludegraphics[width=0.9\linewidth]{artifacts/figures/validity_lag/fig_validity_lag}
\caption{Temporal validity can fail before detector alarm. As drift intensifies,
the validity horizon contracts until the operating horizon becomes
stale at $t_{\mathrm{valid}}$; the detector responds only later at
$t_{\mathrm{detect}}$, leaving a strictly positive validity-detection lag. The
lower panel shows excess tracking error accumulating before the alarm appears.}
\label{fig:validity_lag}
\end{figure}

Across the reported runs, the measured lag remains positive. Retained evidence
can already be invalid for the present while detector statistics remain below
alarm threshold.

The positive lag is not a knife-edge feature of one tuning choice. A compact
false-alarm-calibrated frontier over $H\in\{0.5,0.75,1.0\}$, operating windows
$160$ and $240$, roughness levels $0.01$ and $0.024$, and detector families
ADWIN, Page--Hinkley, KSWIN, and CUSUM shows the same separation on the tested
observation-based detector families. Observation-based ADWIN and Page--Hinkley
retain unit detection rate and strictly positive mean gaps on the successful
grid points, while KSWIN preserves the sign of the gap when it detects but with
substantially lower detection rates. Residualized detector inputs are not
uniformly better on this family. The calibrated sign result therefore holds for
observation-based detector families on the tested compact grid, not as a
detector-uniform delay law; see \Cref{tab:validity_gap_robustness}.

\subsection{Conformal Calibration under Drift}
\label{sec:empirical_conformal}

The same geometry appears when retained information is used for predictive
coverage rather than for direct distribution tracking. In the split-conformal
benchmark, the calibration window acts as a finite-memory state, the
nonconformity quantile is the readout, and the score combines interval width
with coverage loss. Small calibration windows are variance dominated, while
large calibration windows become stale under drift, so the score has an interior
optimum over window size.

The benchmark also separates the calibration-validity clock from the detection
clock. Coverage loss occurs before alarm on the tested drifted residual family,
leaving a positive coverage-before-alarm gap. Around the score-minimizing
window, the empirically safe calibration band coincides with the useful-memory
band. Conformal prediction under drift therefore fits the same validity law:
retained calibration is useful only up to a finite horizon.

\subsection{Empirical Summary}\label{sec:empirical_summary}

The experiments support a structural validity profile with a near-optimal
useful-memory region, roughness-dependent horizon scaling, a corresponding
pattern on the ELEC2 stream, a measurable lag between validity loss and detector
alarm, and a conformal calibration window with the same horizon geometry. The
online adaptation section summarizes the controller and phase-map evidence; it
shows that the structural controller keeps mean error close to the oracle while
the preferred scale rotates with drift phase.
